\section{Introduction}\label{sec:introduction}

Blockchain is a distributed network of computers (\emph{peers}).
Each peer receives transaction requests, that eventually packs
into \emph{blocks}. These form a growing tree,
where each block points back to the previous one. Mining is the process of
block creation and includes a notion of chain quality to
decide which of the many tree branches emerges
as the best chain, on which all peers agree,
hence reaching \emph{consensus}. Mining
is often remunerated: block creators may
mint and earn some amount of cryptocurrency.

This abstract description is explicitly generic \wrt{}
the notion of transaction, irrelevant here: a transaction is an operation
that modifies the state of an abstract machine.
For instance, it is a money transfer in Bitcoin~\cite{Nakamoto08} or
the execution of some code over shared objects in Ethereum~\cite{AntonopoulosW25}.
This description is also generic \wrt{} the mining algorithm.
Bitcoin pioneered proof of work~\cite{DworkN92},
while more modern blockchains mostly use
proof of stake~\cite{Kwon14}, and this paper focuses on proof
of space~\cite{Spoto25}.
In Bitcoin's proof of work,
minted blocks carry evidence of the performance of some computational
work for their creation. The more this work, the higher the
likelihood for the blocks to be integrated in the best chain.
It is well-known that this induces high electricity costs for CPU
computation and data centers, with a high ecological impact (resource
consumption, e-waste). Moreover, this leads to the concentration of mining power
in countries where electricity is cheaper, against the idea of full decentralization.
Proof of stake emerged as a reaction to that. Its idea is that
only a restricted set of peers (\emph{validators}) create blocks.
This makes mining cheaper, efficient and ecologically sustainable.
Non-validators may stake on the success of validators. That introduces some
decentralization, but
proof of stake is generally considered as more centralized than proof of work.
Still, it is the most used mining algorithm.
The idea of proof of space~\cite{AbusalahACKPR17,AtenieseBFG14,CohenP19,DziembowskiFKP15,ParkKFGAP18,RenD16,Spoto25} is that blocks carry evidence
of the allocation of large data on disk, used for their creation.
The larger that data, the higher the likelihood that
blocks get included in the best chain. Proof of space is decentralized like
Bitcoin's proof of work and ecological like proof of stake.

This paper focuses on \emph{mobile mining}, that is,
on mining with a mainstream mobile device (typically, a phone or tablet).
In principle, this is appealing because:
%
\begin{itemize}
\item it leverages a huge existing hardware base, that otherwise remains unused
  for most of the time;
\item it introduces mining to a large portion of humanity, who owns a
  phone but not a powerful connected computer;
\item it can be easily updated, thanks to the automatic update system of mobile devices;
\item it simplifies the start-up of new blockchains, since it is easier to convince people
  to install an app than to run and maintain some software on an always connected computer.
\end{itemize}
%
Mobile mining has not been used yet because of the obvious limitations
of mobile devices \wrt{} battery capacity and network bandwidth, in particular
in least favorite countries.

This paper makes the following contributions:
%
\begin{enumerate}
\item it presents mobile mining, its power and limitations;
\item it shows how proof of space emerges as a niche solution, less sensitive
  to such limitations;
\item it presents Mokaminter, our actual Android app for mobile mining with proof of space;
\item it reports experiments with Mokaminter, discussing its minimal CPU, RAM and data exchange costs
  and its battery use, as well as the amount of crypto earned with different devices.
\end{enumerate}
%
Mokaminter is a miner for mobile devices, not a blockchain
peer: it does not hold nor exchange blocks, it does not host the blockchain,
it does not use a database: all that would be prohibitive on a mobile device.
Mokaminter connects and mines for an external peer that keeps the blockchain.

The rest of this paper is organized as follows.
Sec.~\ref{sec:related_work} discusses related work.
Sec.~\ref{sec:mining} gives a uniform description of various mining algorithms and
discusses if they can be used on a mobile device.
Sec.~\ref{sec:pos} describes Mokamint's proof of space mining algorithm.
Sec.~\ref{sec:mokaminter} describes the Mokaminter mobile app.
Sec.~\ref{sec:experiments} confirms its reduced costs with actual experiments.
Sec.~\ref{sec:conclusion} concludes.
